\section{Role of SE blocks} \label{sec:interpretation}

Although the proposed SE block has been shown to improve network performance on multiple visual tasks, we would also like to understand the relative importance of the squeeze operation and how the excitation mechanism operates in practice. A rigorous theoretical analysis of the representations learned by deep neural networks remains challenging, we therefore take an empirical approach to examining the role played by the SE block with the goal of attaining at least a primitive understanding of its practical function.
%---------------table and figure ----------------------------
\begin{table}[t]
\renewcommand\arraystretch{1.1}
\caption{Effect of Squeeze operator on ImageNet (error rates \%).}
\label{tab:effect_squeeze}
\vspace{-1.7em}
\small
\begin{center}{\scalebox{0.96}{
\begin{tabular}{l|p{1.3cm}<{\centering}|p{1.3cm}<{\centering}|p{1.2cm} <{\centering}|p{1.2cm} <{\centering}}
\hline
& top-1 err. & top-5 err. & GFLOPs & Params\\
\hline
ResNet-50 &   $23.30$ &   $6.55$ & $3.86$ & $25.6$M\\
NoSqueeze & $22.93$ & $6.39$ & $4.27$ & $28.1$M\\
SE & $\mathbf{22.28}$  &  $\mathbf{6.03}$ & $3.87$ & $28.1$M\\
\hline
\end{tabular}}}
\end{center}
\end{table}

\begin{figure*}[t]
\small
\begin{center}
\subfigure[\texttt{SE\_2\_3}]{\includegraphics[trim={1cm 7cm 1cm 7cm},clip,width=0.28\textwidth,height=0.15\textheight]{activations-no-std/conv2_3_sigmoid.pdf}}\hspace{0.2cm}
\subfigure[\texttt{SE\_3\_4}]{\includegraphics[trim={1cm 7cm 1cm 7cm},clip,width=0.28\textwidth,height=0.15\textheight]{activations-no-std/conv3_4_sigmoid.pdf}}
    \hspace{0.2cm}
\subfigure[\texttt{SE\_4\_6}]{\includegraphics[trim={1cm 7cm 1cm 7cm},clip,width=0.28\textwidth,height=0.15\textheight]{activations-no-std/conv4_6_sigmoid.pdf}}
 \hspace{0.2cm} 
\subfigure[\texttt{SE\_5\_1}]{\includegraphics[trim={1cm 7cm 1cm 7cm},clip,width=0.28\textwidth,height=0.15\textheight]{activations-no-std/conv5_1_sigmoid.pdf}}
 \hspace{0.2cm} 
\subfigure[\texttt{SE\_5\_2}]{\includegraphics[trim={1cm 7cm 1cm 7cm},clip,width=0.28\textwidth,height=0.15\textheight]{activations-no-std/conv5_2_sigmoid.pdf}}
 \hspace{0.2cm}
\subfigure[\texttt{SE\_5\_3}]{\includegraphics[trim={1cm 7cm 1cm 7cm},clip,width=0.28\textwidth,height=0.15\textheight]{activations-no-std/conv5_3_sigmoid.pdf}}
\end{center}
\vspace{-0.5cm}
\caption{Activations induced by the Excitation operator at different depths in the SE-ResNet-50 on ImageNet. Each set of activations is named according to the following scheme: \texttt{SE\_stageID\_blockID}.  With the exception of the unusual behaviour at \texttt{SE\_5\_2}, the activations become increasingly class-specific with increasing depth.}
\label{fig:class-activations}
\end{figure*}

\begin{figure*}[t]
\small
\begin{center}
\subfigure[\texttt{SE\_2\_3}]{\includegraphics[trim={1cm 7cm 1cm 7cm},clip,width=0.28\textwidth,height=0.15\textheight]{activations-with-std/conv2_3_sigmoid_std.pdf}} \hspace{0.2cm}
\subfigure[\texttt{SE\_3\_4}]{\includegraphics[trim={1cm 7cm 1cm 7cm},clip,width=0.28\textwidth,height=0.15\textheight]{activations-with-std/conv3_4_sigmoid_std.pdf}} \hspace{0.2cm}
\subfigure[\texttt{SE\_4\_6}]{\includegraphics[trim={1cm 7cm 1cm 7cm},clip,width=0.28\textwidth,height=0.15\textheight]{activations-with-std/conv4_6_sigmoid_std.pdf}}  \hspace{0.2cm}
\subfigure[\texttt{SE\_5\_1}]{\includegraphics[trim={1cm 7cm 1cm 7cm},clip,width=0.28\textwidth,height=0.15\textheight]{activations-with-std/conv5_1_sigmoid_std.pdf}}  \hspace{0.2cm}
\subfigure[\texttt{SE\_5\_2}]{\includegraphics[trim={1cm 7cm 1cm 7cm},clip,width=0.28\textwidth,height=0.15\textheight]{activations-with-std/conv5_2_sigmoid_std.pdf}} \hspace{0.2cm}
\subfigure[\texttt{SE\_5\_3}]{\includegraphics[trim={1cm 7cm 1cm 7cm},clip,width=0.28\textwidth,height=0.15\textheight]{activations-with-std/conv5_3_sigmoid_std.pdf}} \hspace{0.2cm}
\end{center}
\vspace{-1.5em}
\caption{Activations induced by {\it Excitation} in the different modules of SE-ResNet-50 on image samples from the goldfish and plane classes of ImageNet. The module is named  ``\texttt{SE\_stageID\_blockID}''.}
\label{fig:instance-activations}
\end{figure*}

\subsection{Effect of Squeeze}\label{subsection:effect-of-squeeze}
To assess whether the global embedding produced by the squeeze operation plays an important role in performance, we experiment with a variant of the SE block that adds an equal number of parameters, but does not perform global average pooling. Specifically, we remove the pooling operation and replace the two FC layers with corresponding $1\times 1$ convolutions with identical channel dimensions in the excitation operator, namely \textit{NoSqueeze}, where the excitation output maintains the spatial dimensions as input. In contrast to the SE block, these point-wise convolutions can only remap the channels as a function of the output of a local operator. While in practice, the later layers of a deep network will typically possess a (theoretical) global receptive field, global embeddings are no longer directly accessible throughout the network in the \textit{NoSqueeze} variant. The accuracy and computational complexity of both models are compared to a standard ResNet-50 model in Table~\ref{tab:effect_squeeze}. We observe that the use of global information has a significant influence on the model performance, underlining the importance of the squeeze operation.  Moreover, in comparison to the \textit{NoSqueeze} design, the SE block allows this global information to be used in a computationally parsimonious manner.

\subsection{Role of Excitation} \label{subsection:role-of-excitation}
To provide a clearer picture of the function of the excitation operator in SE blocks, in this section we study example activations from the SE-ResNet-50 model and examine their distribution with respect to different classes and different input images at various depths in the network.  In particular, we would like to understand how excitations vary across images of different classes, and across images within a class.

We first consider the distribution of excitations for different classes. Specifically, we sample four classes from the ImageNet dataset that exhibit semantic and appearance diversity, namely \textit{goldfish}, \textit{pug}, \textit{plane} and \textit{cliff} (example images from these classes are shown in Appendix). We then draw fifty samples for each class from the validation set and compute the average activations for fifty uniformly sampled channels in the last SE block of each stage (immediately prior to downsampling) and plot their distribution in Fig.~\ref{fig:class-activations}. For reference, we also plot the distribution of the mean activations across all of the $1000$ classes.

We make the following three observations about the role of the {\it excitation} operation. First, the distribution across different classes is very similar at the earlier layers of the network, e.g. SE\_2\_3.  This suggests that the importance of feature channels is likely to be shared by different classes in the early stages.  The second observation is that at greater depth, the value of each channel becomes much more \mbox{class-specific} as different classes exhibit different preferences to the discriminative value of features, e.g. SE\_4\_6 and SE\_5\_1.  These observations are consistent with findings in previous work \cite{lee_icml2009cdbn, yosinki_nips2014transfer}, namely that earlier layer features are typically more general (e.g. class agnostic in the context of the classification task) while later layer features exhibit greater levels of specificity \cite{morcos_iclr2018importance}.

Next, we observe a somewhat different phenomena in the last stage of the network. SE\_5\_2 exhibits an interesting tendency towards a saturated state in which most of the activations are close to one. At the point at which all activations take the value one, an SE block reduces to the identity operator. At the end of the network in the SE\_5\_3 (which is immediately followed by global pooling prior before classifiers), a similar pattern emerges over different classes, up to a modest change in scale (which could be tuned by the classifiers).  This suggests that SE\_5\_2 and SE\_5\_3 are less important than previous blocks in providing recalibration to the network.  This finding is consistent with the result of the empirical investigation in Section~\ref{sec:modelcapacity} which demonstrated that the additional parameter count could be significantly reduced by removing the SE blocks for the last stage with only a marginal loss of performance.

Finally, we show the mean and standard deviations of the activations for image instances within the same class for two sample classes (\textit{goldfish} and \textit{plane}) in Fig.~\ref{fig:instance-activations}.  We observe a trend consistent with the inter-class visualisation, indicating that the dynamic behaviour of SE blocks varies over both classes and instances within a class. Particularly in the later layers of the network where there is considerable diversity of representation within a single class, the network learns to take advantage of feature recalibration to improve its discriminative performance \cite{hu2018GE}. In summary, SE blocks produce instance-specific responses which nevertheless function to support the increasingly class-specific needs of the model at different layers in the architecture.